\section{Introduction}
\label{sec:introduction}

Clinical diagnosis is inherently multimodal. When a gastroenterologist reads an
endoscopic image, the judgment fuses two information streams: visual patterns
(mucosal texture, lesion morphology, staining response) and textual knowledge
(clinical guidelines, prior case descriptions). Automated medical question
answering (QA) has begun to mirror this process, moving from closed-book
vision-language models~\cite{li2023llavamed} toward retrieval-augmented
generation (RAG)~\cite{lewis2020retrieval}, where external evidence is retrieved
to ground the answer. In the multimodal setting the evidence is itself
image--text: an ideal system retrieves relevant images \emph{and} text passages,
then reasons over them jointly.

Yet a blind spot pervades current multimodal medical RAG. Although images are
retrieved, most pipelines convert them into textual
surrogates---captions or descriptions---and feed only text to the
generator~\cite{chen2022murag,yasunaga2023retrieval}; the generation model never
actually sees the retrieved image. On EndoBench~\cite{endobench}, a large
endoscopic QA benchmark, $84.5\%$ of questions are image-dependent. For these
questions a text-proxy pipeline systematically hides the value of visual
evidence, so any gain from better visual retrieval or reranking becomes
invisible under evaluation. We argue that once the actual retrieved images are
fed to the generator, the bottleneck of multimodal medical RAG shifts from
recall to \emph{cross-modal evidence selection}: choosing, from a mixed pool of
image and text candidates, the small set that best supports the answer.
Figure~\ref{fig:motivation} illustrates this gap.

This shift exposes three problems that existing methods do not address.

\paragraph{(P1) Selection is passive, and relevance is not answer-utility.}
Rerankers score each candidate for query relevance and take the top $K$ in a
single pass. This has two consequences. First, relevance is only a proxy: a
passage can be topically on-target yet contribute nothing to---or even distract
from---answering the specific question, and a similarity score cannot separate
the two. What ultimately matters is whether a piece of evidence raises the
generator's chance of producing the correct answer. Second, the pass is
one-shot: when the retrieved pool is poor, the system has no recourse. It cannot
reformulate the query and search again, but must answer from whatever the first
retrieval returned.

\paragraph{(P2) In medicine, visual similarity does not imply clinical identity.}
Endoscopic images are visually homogeneous: distinct pathologies at the same
anatomical site appear nearly identical in pixel space. Pure similarity-based
retrieval and top-$K$ selection consequently return confident but clinically
wrong evidence---a failure mode largely absent from generic RAG. Effective
selection must judge whether a candidate actually \emph{discriminates} the
answer, not merely whether it looks similar.

\paragraph{(P3) Low-discriminability text is a noise source, not merely a weak signal.}
The discriminative content of a query is distributed unevenly across
modalities, and, more sharply, some query text actively harms retrieval.
Generic or positional phrasings---``which anatomical site is shown,'' ``what
structure is this''---carry no discriminative retrieval signal; issued as a text
query they pull in candidates that share the same site or structure but differ
in pathology, systematically injecting homogeneous false positives. The problem
is thus not only that a fixed fusion weight mis-balances the two paths, but that
a modality can be worth \emph{suppressing}: retrieving more text can lower answer
quality. Effective fusion must decide, per query, how much to trust each
modality---and when to gate one off entirely.

We present \textbf{MedAlign-RAG}, a multimodal medical RAG framework built on a
single insight: recast fusion-stage evidence selection from maximizing
query--candidate relevance into an \emph{agentic} decision process trained by
downstream answer-utility. A vision-language agent inspects the reranked pool,
judges the value of each image and text candidate in the generator's own visual
space, and chooses between two actions: accept the evidence and generate, or
rewrite the query and re-retrieve. The agent is trained by reinforcement
learning whose reward is the increase in the (frozen) generator's probability of
the correct answer once the selected evidence is supplied---a signal available
directly from multiple-choice supervision, requiring no passage-level
annotation. A density-aware modality router allocates the text/image retrieval
budget on every round and suppresses a modality whose signal is noise.

Our contributions are:
\begin{itemize}
  \item We recast multimodal evidence selection as an \emph{agentic search}
  process and train a vision-language agent with an \emph{answer-utility
  reward}---the lift in the generator's probability of the correct answer---so
  evidence is optimized for its effect on the answer rather than for surface
  relevance, using only multiple-choice supervision and no passage-level labels.
  \item The agent operates \emph{cross-modally}, judging image and text
  candidates alike and, when the pool is inadequate, rewriting the query to
  re-retrieve; this recovers from poor first-stage retrieval and directly targets
  the visually-homogeneous false positives characteristic of medical images.
  \item We propose \emph{density-aware modality routing}, which predicts a
  query's textual density and image dependency to allocate---and, when a modality
  is noise, suppress---the retrieval budget across the text and image paths on
  every retrieval round.
  \item We identify an \emph{evaluation blind spot} in multimodal medical
  RAG---retrieving images but generating from text alone---and evaluate
  MedAlign-RAG on EndoBench under a corrected end-to-end image--text protocol,
  comparing against no-retrieval, top-$K$, cross-encoder, graph-based, and
  agentic-RAG baselines.
\end{itemize}

\begin{figure}[t]
  \centering
  \IfFileExists{figures/fig1_motivation.pdf}{%
    \includegraphics[width=\columnwidth]{fig1_motivation.pdf}}{%
    \framebox[0.95\columnwidth]{\rule{0pt}{3.4cm}\ Figure~1 placeholder (motivation)}}
  \caption{Motivation. Existing multimodal medical RAG retrieves images but
  passes only their textual descriptions to the generator (text-proxy). Because
  most medical queries are image-dependent, this hides the value of visual
  evidence; feeding the actual retrieved images exposes cross-modal evidence
  selection as the true bottleneck.}
  \label{fig:motivation}
\end{figure}
